\documentclass[12pt]{article}

\title{Journal for Complex Analysis}
\author{Michael Gavina}


\usepackage{amssymb,amsfonts,amsmath,amsthm,url,graphicx}
\graphicspath{ {./images/} }


%\pagestyle{empty} % for no page numbers
\setlength{\oddsidemargin}{-.1in}      % 1in left margin
\setlength{\evensidemargin}{-.1in}     % 1in left margin (even pages)
\setlength{\topmargin}{0in}           % 1.50in top margin
\setlength{\textwidth}{6.7in}           % 6.3in text - 1.25in rt margin
\setlength{\textheight}{9in}            % Body ht for 1in margins
\addtolength{\topmargin}{-\headheight}  % No header, so compensate
\addtolength{\topmargin}{-\headsep}     % for header height and separation
%  End of margins and formatting params %
\linespread{1}
\setlength{\parindent}{0pt}


\usepackage{color}
\definecolor{cherry}{rgb}{0.9,.1,.2}
\definecolor{green}{rgb}{.2,.7,0}
\definecolor{blue}{rgb}{0,.0, .81}

\def\green#1{\textcolor{green}{#1}}
\def\cherry#1{\textcolor{cherry}{#1}}
\def\blue#1{\textcolor{blue}{#1}}

\begin{document}
\maketitle

\begin{center}
\textbf{\large Chapter 0}
\end{center}

\noindent{\bf Task 1}
\\
\noindent{\bf a}
$x \mapsto 2x$ \\
Here, the original function is transformed in which only the positive and negative even numbers and zero are shown on the number line. This transformation is similiar to a dialation of a geometric shape.
\\

\noindent{\bf b}
$x \mapsto 2x + 1$ \\
Here, the original function is transformed in which only the the positive and negative odd numbers are shown on the number line. This transformation first dilates then translates the line by one.
\\

\noindent{\bf c}
$x \mapsto -x$ \\
Here, the original function is transformed into the same number line as before. However, the effect the transformation is that of a reflection at zero. This set is an image of the original function. 
\\

\noindent{\bf d}
$x \mapsto -x + 4$ \\
Here, the original function is transformed by a reflection on the line at zero then a translation will be applied to shift the line by four. However, this new set is an image of the original function. 
\\

\noindent{\bf e}
$x \mapsto -3x + 6$ \\
Here, the original function is transformed by a reflection then a translation across the line. Within this new function, the only numbers contained within this set, would be powers of three.
\\

\noindent{\bf Task 2}\\
\noindent{\bf a}
$x \mapsto x^2$ \\
Here, the original function is dilated more as the numbers increase, however, this function only allows positive outputs from the intial values. This means that only half of the number line would be used when the transformation is applied.
\\

\noindent{\bf b}
$x \mapsto x^3$ \\
Here, the original function is similiarly dialated to the last function, however, this time, negative numbers can be a possible output of the function. This allows the full number line to be used when representing the function.
\\

\noindent{\bf c} 
$x \mapsto \frac{1}{x}$ \\
Here, the transformation of this function makes its so that all of the numbers are contained between 0 and 1, and 0 and -1. However, zero is not included within this interval because of its undefined behavior. \\
\noindent{\bf Task 3} \\

\noindent{\bf a} 
$x \mapsto 10^x$ \\
Here, the transformation of the function stretches out the positive side of the line rather extremely while the negative side is non-exsitent. In fact, the line can never cross zero at all, leaving zero as a hole within the line. \\

\noindent{\bf b}
$x \mapsto e^x$ \\
This transformation is similar to the last transformation, however, the factor in which the postive side of the line will stretch will be weaker than the last transformation.\\

\noindent{\bf c} 
$x \mapsto e^{-x}$ \\
This transformation is different from the previous two; the negative numbers will produce postive values which will increase the more negative the number is, while the postive values will only produce values closer and closer to zero, and zero is still a hole.\\

\noindent{\bf d}  
$x \mapsto -e^x$ \\
This transformation will produce the oppisite of the $e^x$ transformation, in which the negative side of the number line will contain numbers while the positve side cannot contain numbers.
This transformation is like a reflection across zero when compared to the $e^x$ transformation. \\

\noindent{\bf e} 
$x \mapsto -e^{-x}$ \\
This transformation is similiar to the last transformation, in which, it too, is like a reflection across zero, when compared to the $-e^x$ function.\\

\noindent{\bf f} 
$x \mapsto \sin(x)$ \\
Here this transformation is different from the other transformation because when the transformation is finished, the domain is restricted to -1 to 1. 
This is because sine is restricted to $-\frac{pi}{2}$ to $\frac{pi}{2}$, any values greater than the domain, the range will be previous values. \\

\noindent{\bf g} 
$x \mapsto \cos(x)$ \\
This transformation is very similiar to the last transformation, however, cosine has a different domain. 
The domain for cosine is $0$ to $pi$, after this, any domain values will produce previous range values.\\

\noindent{\bf Task 4}\\
There are two parts to this question.\\ 

Why are properties (F5) and (F6) not true for $\mathbb{N}$?\\

For F5, given that natural numbers are all of the numbers that are not negative, there is not an addivitive inverse of any of the natural numbers included within the natural numbers set. \\

For F6, given that natural numbers are all positve whole numbers, there is never a number in which $\frac{p}{q}$ can ever be true. Thus, there is never a multiplicative inverse.\\


\noindent{\bf Task 5}\\
There are two parts to this question.\\

Which properties are true for each set of sets: $\mathbb{Z}$, $\mathbb{Q}$?\\

For $\mathbb{Z}$, every property will apply except for F6, in which, there is never such a case in which a fractional number can be apart of this set. 
Since the numbers are all integers, the basic rules apply for all intergers (F1-F3), Here $0$ and $1$ are included within this set (F4), and negative numbers exist within this set (F5).\\

For $\mathbb{Q}$, every property will apply for this set. 
Since the numbers are fractional, they can be added and multiplied. 
This also allows for the communtaive and associative property to be true as well and thus (F1-F3) are true.
$0$ and $1$ are included within this set and thus (F4) is true. 
Negative fractions and positive fractions of the same value can exist within this set and thus (F5) is true.
Inverse fractions of fractional and whole numbers do exist within this set and thus (F6) is true.\\

\noindent{\bf Task 6}\\

The rules of ordered sets are that of which there are numbers within the set that are bigger and that adding two numbers will create a bigger number than $0$ within the set.
Here, it is possible to have fractional elements that will exceptionally small numbers but nonetheless still be bigger than $0$. 
The only way in which $0$ could be a possible would to have $0$ as one of the possible numbers as $a$ or $b$, which breaks the ruling for order sets.
Thus, $\mathbb{Q}$ is an ordered set.\\

\noindent{\bf Task 7}\\

Within A, the number $\frac{1}{2}$ resides, while within B, the number $2$ is within the set.\\

Within both sets, there are restrictions on the value of both $q$ and $p$. 
Solving the inequality shows that for set A, the fractional representation that is outputed from the set must follow this ruling: $\frac{p^2}{q^2} < 2$.
Similiarly, there a inequality that must be followed for set B: $\frac{p^2}{q^2} > 2$. 
Here it is clearly shown that set B will contain bigger numbers always, because it's fractional representation must always be bigger than 2.\\

However, both of these sets are not Dedekind complete because of the rulings. 
Both sets can never contain the irrational number $\sqrt(2)$. 
Thus, in both sets, there is a hole located at $\sqrt(2)$.\\

\noindent{\bf Task 8}\\

For $\mathbb{N}$, this is representation of the number line only containing, positive whole numbers aswell as $0$.\\

For $\mathbb{Z}$, this is representation of all of the positive and negative whole numbers along with $0$.\\

For $\mathbb{Q}$, this is representation of all of the numbers across the number line that are not irrational.\\

\noindent{\bf Task 9}\\

\begin{center}
	\begin{tabular}{| c c c |}
	\hline
		Task 1 & ]0,1[ & [2,5]\\
	\hline\hline
		A. & ]0,2[ & [4,10]\\
	\hline
		B. & ]0,3[ & [5,11]\\
	\hline 
		C. & ]-1,0[ & [-5,-2]\\
	\hline
		D. & ]3,4[ & [1,2]\\
	\hline
		E. & ]3,6[ & [-11,0]\\
	\hline
	\end{tabular}
\end{center}


\begin{center}
	\begin{tabular}{| c c c |}
	\hline
		Task 2 & ]0,1[ & [2,5]\\
	\hline\hline
		A. & ]0,1[ & [4,25]\\
	\hline
		B. & ]0,1[ & [8,125]\\
	\hline 
		C. & ]0,1[ & [$\frac{1}{5}$,$\frac{1}{2}$]\\
	\hline
	\end{tabular}
\end{center}


\begin{center}
	\begin{tabular}{| c c c |}
	\hline
		Task 3 & ]0,1[ & [2,5]\\
	\hline\hline
		A. & ]0,10[ & [100,100000]\\
	\hline
		B. & ]1,$e$[ & [$e^2$, $e^5$]\\
	\hline 
		C. & ]0,$\frac{1}{e}$[ & [$\frac{1}{e^5}$, $\frac{1}{e^2}$]\\
	\hline
		D. & ]$-e$, 1[ & [$-e^5$, $-e^2$]\\
	\hline
		E. & ]$-\frac{1}{e}$,0[ & [$-\frac{1}{e^2}$,$-\frac{1}{e^5}$]\\
	\hline
		F. & ]0,$sin(1)$[ & [$-1$, $sin(2)$]\\
	\hline
		G. & ]$cos(1)$,1[ & [$-1$, $cos(5)$]\\
	\hline
	\end{tabular}
\end{center}

\end{document}
